Double descent (DD) predicts that test risk peaks at the interpolation threshold and falls again in the over-parameterised regime. We reproduce textbook DD on random Fourier features (RFF) for MNIST, mechanistically attribute the peak to feature-Gram conditioning collapse via ridge and noise sweeps, and confirm a bias--variance decomposition with variance dominating the peak. The textbook neural-network DD curve is harder to surface: a naive CNN/ResNet width sweep on $n=4{,}000$ noisy CIFAR-10 fails because raw parameter count overshoots the effective interpolation threshold. We repair the capacity axis with a \emph{fractional-$k$} ResNet whose stage widths scale continuously through the threshold; on this family DD reappears as a clean rise--peak--valley--recovery curve, with final-epoch test accuracy moving $26\%\!\to\!49\%\!\to\!53\%$ across $k\in[0.0625, 2]$ at $15\%$ label noise. Five independent diagnostics --- penultimate-feature stable rank, full empirical NTK, Bartlett (2020) effective rank, sample-wise peak migration, and Hessian top eigenvalue --- localise the same phase transition near $k\approx0.1875$. Finally we audit the evaluation protocol: reporting the best-over-epochs test accuracy systematically shallows the over-parameterised valley relative to final-epoch reporting, with the bias concentrated in exactly the regime where DD claims most depend on it.
